\section{The primitive two-loop three-point graph}
\label{sec:primitive}

This chapter treats L2E3G1, the graph $K_{2,3}$ with one fixed-label external
leg at the midpoint of each of its three length-two branches. Use the common
mass and normalized measures of Section~\ref{sec:scheme}. All momenta are
real Euclidean and incoming, with $p_1+p_2+p_3=0$ and $s_i=p_i^2$.
The formulas cover every such three-point configuration, including
degenerate Gram matrices and the zero point.

\subsection{Routing and subtraction}

Choose branch flows $t_1=k$, $t_2=q$, $t_3=-k-q$. The integrand is
\[
 I(p_1,p_2)=D(k)D(k+p_1)D(q)D(q+p_2)
             D(k+q)D(k+q-p_3).
\]
Every proper loop-containing 1PI subgraph has at least four internal edges;
its one-loop degree is at most $6-8=-2$. There are no proper UV-divergent
subgraphs, whereas $L=2$, $I=6$ gives overall degree zero. Therefore
\[
 \overline R I=I,\qquad RI=I(p_1,p_2)-I(0,0).
\]
The scalar answer below does not include the graph weight $1/S_G=1/2$.

\subsection{Gaussian minimization and the exact parameter polynomials}

Let the two parameters on branch $i$ be $a_i,b_i>0$, and put
$A_i=a_i+b_i$, $r_i=b_i/A_i$. Then
\[
 a_it_i^2+b_i(t_i+p_i)^2
 =A_i(t_i+r_ip_i)^2+A_ir_i(1-r_i)s_i.
\]
The constraint $\sum_i t_i=0$ makes the sum of the completed-square
vectors equal to $\sum_i r_ip_i$. Minimizing their weighted squared norm
gives the extra cost
\[
 \frac{|\sum_i r_ip_i|^2}{\sum_i A_i^{-1}}.
\]
Write $(A_1,A_2,A_3)=(A,B,C)$ and $(r_1,r_2,r_3)=(u,v,w)$. The first
Symanzik polynomial and the effective external scale are
\begin{align*}
 U&=AB+AC+BC,\\
 \Omega_p&=Au(1-u)s_1+Bv(1-v)s_2+Cw(1-w)s_3\\
 &\quad+\frac{ABC}{U}|(u-w)p_1+(v-w)p_2|^2,\\
 F&=U\{M(A+B+C)+\Omega_p\}.
\end{align*}
Here $\Omega_p$ denotes an external scale, not the superficial UV degree.
The final square can equally be written
\[
 s_1(u-w)^2+s_2(v-w)^2
 +(s_3-s_1-s_2)(u-w)(v-w).
\]
It is nonnegative exactly on the physical Gram domain. These expressions
are homogeneous before imposing any projective normalization. An independent
spanning-tree and two-forest construction of $U,F$ can therefore be checked
against them without using the same Gaussian derivation.

\subsection{Common cutoff and the finite projective integral}

Apply the same lower cutoff $\alpha_e\ge\eta$ to all six original edge
parameters and subtract the zero-momentum regulated graph first. Set
$\alpha_e=\rho x_e$, $\sum_e x_e=1$. The six-parameter Jacobian contributes
$\rho^5$, while the two six-dimensional Gaussian integrals contribute
$\rho^{-6}$. The remaining radial measure is $d\rho/\rho$, with the
same lower limit $T(x)=\eta/\min_e x_e$ for both terms.
Define
\[
 L_T(z)=\int_T^\infty\frac{d\rho}{\rho}
                  e^{-M\rho}(1-e^{-z\rho}).
\]
For $z\ge0$,
\[
 0\le L_T(z)\le\log(1+z/M),\qquad
 L_T(z)\longrightarrow\log(1+z/M).
\]
Within a branch, $da_i\,db_i=A_i\,dA_i\,dr_i$. On $A+B+C=1$ the full
projective measure is consequently
$ABC\,dA\,dB\,du\,dv\,dw$. Thus the candidate finite formula is
\begin{align*}
 J_{\mathrm{pr}}(s_1,s_2,s_3;M)
 =-&\int_{\substack{A,B\ge0\\A+B\le1}}dA\,dB
       \int_{[0,1]^3}du\,dv\,dw\\
 &\times\frac{ABC}{(AB+AC+BC)^3}
           \log\left(1+\frac{\Omega_p}{M}\right),
 \qquad C=1-A-B.
\end{align*}
The density is nonpositive; its parameter $L^1$ norm equals
$-J_{\mathrm{pr}}$. Its finiteness and the regulator limit follow from
the next estimates, rather than from separately integrating divergent terms.

For comparison only, in bare dimensional regularization with
$d=6-2\epsilon$, the radial Gamma factor is $\Gamma(2\epsilon)$. The
projective period computed below is $1/2$, giving the unweighted simple
pole $1/(4\epsilon)$. There is no extra factor of $1/2$ multiplying the
finite BPH logarithm, since
\[
 \Gamma(2\epsilon)
 \bigl(F^{-2\epsilon}-F_0^{-2\epsilon}\bigr)
 \longrightarrow-\log(F/F_0).
\]
The finite Taylor scheme is not replaced by the bare Laurent finite part.

\subsection{Exact period, compact control, and global scale behavior}

The trial branch currents $t_i=-p_i/2$ obey $\sum_i t_i=0$. The Gaussian
minimum cannot exceed their value. On $A+B+C=1$ this proves
\[
 0\le\Omega_p\le\frac{As_1+Bs_2+Cs_3}{4}
                    \le\frac{s_{\max}}4,
 \qquad s_{\max}=\max_i s_i.
\]
The weight $ABC/U^3$ is homogeneous of degree $-3$. Put
$A=x/(1+x+y)$, $B=y/(1+x+y)$, $C=1/(1+x+y)$. Its projective Jacobian is
$(1+x+y)^{-3}$, so its period is exactly
\begin{align*}
 \mathcal P
 &=\int_{A+B+C=1}\frac{ABC}{(AB+AC+BC)^3}\,dA\,dB\\
 &=\int_0^\infty\int_0^\infty
             \frac{xy\,dx\,dy}{(xy+x+y)^3}\\
 &=\int_0^\infty\frac{dy}{2(1+y)^2}=\frac12.
\end{align*}
The successive integrations are justified by Tonelli's theorem. They prove
integrability at every parameter face as well as fixing the normalization.
It follows that
\[
 |J_{\mathrm{pr}}|=\int|\text{parameter density}|
 \le\frac12\log\left(1+\frac{s_{\max}}{4M}\right)
 \le\frac{s_{\max}}{8M}.
\]
The same integrable weight times this constant controls $L_T$, uniformly
in $\eta$. Parameter dominated convergence therefore removes the common
cutoff. If $|p_i|\le P$ and $M\ge m_0^2>0$, the common bound is
\[
 |J_{\mathrm{pr}}|\le
 \frac12\log\left(1+\frac{P^2}{4m_0^2}\right).
\]
Unlike a fixed-point finiteness assertion, the logarithmic estimate also
controls arbitrarily large common external scales.

Along the ray $s_i\mapsto\lambda s_i$, $\Omega_p$ is linear in $\lambda$.
At $\lambda=1$, differentiation under the positive integrable weight gives
\[
 \left|\frac{d^nJ_{\mathrm{pr}}}{d\lambda^n}\right|
 \le\frac{(n-1)!}{2}
       \min\left\{\left(\frac{s_{\max}}{4M}\right)^n,1\right\},
 \qquad n\ge1.
\]
The first derivative is nonpositive and the second nonnegative. If the
shape is fixed and $s_{\max}$ itself is the scale variable, the bound is
$\frac12(n-1)!\min\{(4M)^{-n},s_{\max}^{-n}\}$ away from zero,
with the continuous small-scale bound at zero.

For independent vectors $z=(p_1,p_2)$, write $\Omega_p=z^TQz$. The trial
bound implies $Q\ge0$ and $\|Q\|\le1/2$, hence
\[
 |\nabla\Omega_p|^2\le2\Omega_p,\qquad
 \|\nabla^2\Omega_p\|\le1.
\]
The period is $1/2$, and elementary maximization of the differentiated
logarithm gives, with $P^2=s_{\max}$,
\begin{align*}
 \|\nabla_{(p_1,p_2)}J_{\mathrm{pr}}\|
 &\le\min\left\{\frac{P}{\sqrt2 M},
                         \frac1{2\sqrt{2M}}\right\},\\
 \|\nabla^2_{(p_1,p_2)}J_{\mathrm{pr}}\|
 &\le\frac3{4M}.
\end{align*}
These are analytic amplitude and derivative bounds, not discretization
error estimates.

\subsection{Independent absolute convergence in the original twelve dimensions}

Parameter absolute integrability after Gaussian integration alone does not
prove the original momentum $L^1$ statement. We now prove the latter.
Put $P=\max_i|p_i|$, $\Lambda=1+P/m$. The first directional derivative
bound is $|\partial D|\le2\langle r\rangle^{-3}$. Along the physical ray,
\[
 RI=\int_0^1\frac{d}{d\tau}I(\tau p_1,\tau p_2)\,d\tau.
\]
There are three terms, one for the shifted propagator on each branch.
All estimates below are uniform in $\tau\in[0,1]$.

\paragraph{One soft branch.}
Consider $|k|\ge2(|q|+P+m)$. Both lines on the $k$ and $k+q$ branches are
hard and comparable with $\langle k\rangle$. Derivatives on the soft branch
are bounded by
$512P\Lambda^3\langle k\rangle^{-8}\langle q\rangle^{-5}$;
the sum of the two hard-branch derivatives is bounded by
$2048P\Lambda^2\langle k\rangle^{-9}\langle q\rangle^{-4}$.
Since $\langle q\rangle\le\langle k\rangle$, both are controlled by
\[
 |RI|\le P C_h(\Lambda)
       \langle k\rangle^{-8}\langle q\rangle^{-5},\qquad
 C_h(\Lambda)=512\Lambda^3+2048\Lambda^2.
\]
The normalized angular factor in a six-dimensional radial integral is one.
Consequently
\begin{align*}
 \int_{|k|\ge2(|q|+P+m)}\langle k\rangle^{-8}\frac{d^6k}{\pi^3}
 &\le\frac1{8\langle q\rangle^2},\\
 \int_{\mathbb R^6}\langle q\rangle^{-7}\frac{d^6q}{\pi^3}
 &=\frac8{15m}.
\end{align*}
Permuting the three branch flows, or making a loop-variable linear change
with absolute determinant one, gives three such soft-branch regions.
Their union has absolute integral at most $PC_h/(5m)$, even if they overlap.
One of them is the soft $k+q$ region, so the channel $k\simeq-q$ is
explicitly included.

\paragraph{Comparable large branch flows.}
Let $K=\max(|k|,|q|,|k+q|)$ and $K_0=8(P+m)$. Outside the soft-branch
hierarchies, if $K\ge K_0$, all three base flows are at least $K/4$;
their physical shifts are at least $K/8$. Therefore
\[
 |RI|\le6P\,8^{13}K^{-13}.
\]
For a dyadic shell $K\in[K_n,2K_n]$, $K_n=2^nK_0$, the product of the
$k,q$ ball volumes is at most $(1024/9)K_n^{12}$. The shell integral is
bounded by a constant times $P/K_n$. Summing
$\sum_{n\ge0}K_n^{-1}=2/K_0$ gives the explicit bound
\[
 \frac{4096}{3}\frac{P\,8^{13}}{K_0}.
\]
This argument does not assume all six line momenta are independently large.

\paragraph{Compact remainder.}
For $K\le K_0$, positive mass bounds the derivative sum by $6Pm^{-13}$.
The product ball volume is at most $K_0^{12}/36$. Thus, for example,
\[
 \int_k\int_q |RI|
 \le\frac{PC_h(\Lambda)}{5m}
       +\frac{4096P\,8^{13}}{3K_0}
       +\frac{P K_0^{12}}{6m^{13}}<\infty.
\]
This deliberately coarse momentum norm differs from the much smaller
post-Gaussian parameter norm. Its controls are uniform when $P$ is bounded
and $m\ge m_0>0$.

The same propagator and first-derivative estimates hold for $D_\eta$, by
the positive Schwinger moment argument in Section~\ref{sec:scheme}. The
three regions therefore provide a uniform momentum majorant and prove
$\int RI_\eta\to\int RI$. Parameter dominated convergence identifies this
same left side with the logarithmic formula above. This completes the
equivalence to the ordinary absolutely convergent BPH momentum integral.

\subsection{A finite unit-cube parameterization}

Partition the branch-sum simplex according to the largest branch and the
order of the other two. On one of the six ordered sectors take
\[
 A=\frac{rt}{1+r+rt},\qquad
 B=\frac r{1+r+rt},\qquad
 C=\frac1{1+r+rt},\qquad 0\le r,t\le1.
\]
The exact measure transformation is
\[
 \frac{ABC}{U^3}\,dA\,dB
       =\frac{t}{(1+t+rt)^3}\,dr\,dt.
\]
Together with $u,v,w\in[0,1]$, this gives six finite five-dimensional
unit-cube integrals. The remaining reduced denominator is at least one,
and the logarithm is bounded by the already proved physical scale bound.
The six invariant permutations must be summed; equal permutations may be
combined by multiplicity. Any subsequent numerical refinement diagnostic
remains distinct from the analytic claims established here.
